This chapter reviews the work most relevant to a defense-oriented risk scorer
for the Model Context Protocol (MCP). We cover five areas: the MCP attack
surface, MCP-specific defenses and attacks, the safety of tool-using agents,
guardrails and runtime gating, and risk-scoring methodologies from CVSS to
LLM-based judges. We then position the proposed research.

\section{The MCP Attack Surface}
\label{sec:related_work:mcp_surface}
\addcontentsline{tocheb}{section}{\protect\numberline{\secnumforhebrewtoc}{משטח התקיפה של MCP}}

MCP standardizes how LLM agents discover and invoke external tools, resources,
and prompts exposed by a server~\cite{anthropic2024mcp}. The uniform interface
that makes MCP powerful also creates a broad attack surface: any agent that
speaks the protocol can request any exposed operation, and the server has little
native machinery to decide whether a request should be allowed --- a gap the
protocol's own security guidance acknowledges~\cite{mcp2026securitybp}.
\citet{hou2025mcp} give the first systematic map of this surface, defining the
server lifecycle and a threat taxonomy across four attacker types, while
\citet{rostamzadeh2026mcpdpt} organize attacks by the architectural component
responsible for enforcement and find protection uneven and predominantly
tool-centric --- directly motivating a server-side gate. Empirical studies
confirm the risk is real: \citet{radosevich2025mcpsafety} coerce models through
MCP tools into compromising a host; \citet{hasan2025mcp} find that of $1{,}899$
servers, $7.2\%$ carry general vulnerabilities and $5.5\%$ exhibit tool
poisoning; \citet{zhao2026parasites} measure $12{,}230$ tools across $1{,}360$
servers and define parasitic toolchain attacks; and \citet{li2025privilege}
statically analyze $2{,}562$ servers and argue for explicit privilege
management. \citet{yang2025mcpsecbench} standardize this into a benchmark of
seventeen attack types and report that existing protections block fewer than
$30\%$ of attacks.

\section{MCP-Specific Defenses and Attacks}
\label{sec:related_work:mcp_defenses}
\addcontentsline{tocheb}{section}{\protect\numberline{\secnumforhebrewtoc}{הגנות והתקפות ייחודיות ל-MCP}}

On the defensive side, \citet{xing2025mcpguard} propose \emph{MCP-Guard}, a
proxy-based three-stage pipeline (static scan, neural detector, and LLM
arbitrator) whose layered architecture is the closest in spirit to ours, and
\citet{jing2025mcip} track contextual integrity of information flows. Because a
tool's natural-language description drives agent tool selection, that
description is itself an attack vector: \citet{shi2025toolhijacker} show a single
poisoned tool document can hijack tool selection, and tool-squatting and
rug-pull attacks motivate OAuth-signed, versioned tool
definitions~\cite{bhatt2025etdi}. The work closest to ours places enforcement at
the server boundary: \citet{buhler2026agentbound} present the first
access-control framework for MCP servers, \citet{sharma2026pauth} propose precise
task-scoped authorization in which the server verifies a call is a legitimate
consequence of the user's task, and \citet{huang2026caller} measure ``caller
identity confusion'' --- servers generally cannot tell which agent is invoking a
request. All three issue an allow/deny decision; our proposal instead computes a
graduated risk \emph{score}.

\section{Safety of Tool-Using Agents and Prompt Injection}
\label{sec:related_work:agent_safety}
\addcontentsline{tocheb}{section}{\protect\numberline{\secnumforhebrewtoc}{בטיחות סוכנים והזרקת פקודות}}

Because the threat originates from an agent, the agent-safety literature
characterizes what a server-side scorer must anticipate.
\citet{yuan2024rjudge} show that even strong models judge the safety of agent
actions correctly only $74\%$ of the time; \citet{andriushchenko2025agentharm}
find leading models compliant with malicious agent tasks; and
\citet{ruan2024toolemu} use LM-emulated sandboxing to measure risky behavior at
scale. Most pointedly for us, \citet{cartagena2026gap} demonstrate that textual
safety does \emph{not} transfer to tool-call safety --- models often refuse in
chat yet still execute the harmful call --- a direct argument for scoring the
call itself. Indirect prompt injection is the mechanism that turns a benign
agent into a threat: \citet{zhan2024injecagent} and \citet{debenedetti2024agentdojo}
show that tool-returned data can hijack agent behavior in realistic workflows.
Rather than detecting injection, our proposal bounds the \emph{damage} an
injected agent can do by scoring each requested operation before it runs.

\section{Guardrails and Runtime Gating}
\label{sec:related_work:guardrails}
\addcontentsline{tocheb}{section}{\protect\numberline{\secnumforhebrewtoc}{מעקות בטיחות ושערים בזמן ריצה}}

A growing line of work places an interpretable enforcement layer between a model
and its effects, from LLM-based input/output classifiers~\cite{inan2023llamaguard}
to per-call least-privilege systems: Progent expresses tool-call permissions as
verifiable symbolic rules~\cite{shi2025progent}, and AgenTRIM reconstructs tool
interfaces offline and blocks misaligned calls at
runtime~\cite{betser2026agentrim}. These systems decide \emph{whether} to allow
an action; our framework is complementary in computing \emph{how risky} an action
is on a graduated scale, a signal such gates could consume. The need is
underscored by evidence that agents can behave as strategic ``insider threats''
when their goals are obstructed~\cite{anthropic2025misalignment}.

\section{Risk Scoring: From CVSS to LLM Judges}
\label{sec:related_work:scoring}
\addcontentsline{tocheb}{section}{\protect\numberline{\secnumforhebrewtoc}{ניקוד סיכון: מ-CVSS ועד שופטי LLM}}

Our framework is, at its core, a scoring system. \citet{spring2021time} argue
CVSS is inadequate as a risk proxy because its algorithm is neither formally nor
empirically justified; \citet{jacobs2021epss} complement severity with
data-driven exploitation likelihood; and \citet{bahar2024validity} show CVSS,
DREAD, and OWASP Risk Rating produce minimal variation across very different LLM
attacks. The community has responded with agentic-specific scores and
checklists~\cite{owasp2025aivss,mehra2025graphofeffort,owasp2026agentictop10},
and with LLMs applied directly to security scoring: prior work from our own group
uses LLMs to detect and reason about misconfigurations
(GenKubeSec~\cite{malul2024genkubesec}, KubeGuard~\cite{cohen2025kubeguard}), and
others evaluate LLMs for vulnerability triage~\cite{haddad2025prompting}. The
mechanism we adopt --- an LLM producing a numeric score against an explicit
rubric --- descends from the LLM-as-judge literature: G-Eval couples
chain-of-thought with form-filling for calibrated scores~\cite{liu2023geval},
and \citet{zheng2023judging} both validate the approach and catalog its biases.
Two very recent works retarget this to our exact setting --- SkillVetBench scores
agent skills with a CVSS-flavored rubric~\cite{hossain2026skillvet}, and
\citet{fleming2025tbac} have an LLM judge each access request to drive
risk-adaptive access control --- and, like them, we validate the LLM's scores
against external oracles rather than trusting them blindly. Finally, since an MCP
call is structurally an authenticated API request, the runtime half of our
framework draws on dynamic risk assessment~\cite{cheimonidis2023dynamic}.

\section{Positioning of the Proposed Research}
\label{sec:related_work:positioning}
\addcontentsline{tocheb}{section}{\protect\numberline{\secnumforhebrewtoc}{מיקום המחקר המוצע}}

Three gaps recur across this literature. First, the fast-growing set of MCP
defenses is overwhelmingly \emph{binary} (attack vs.\ benign) and tool-centric,
leaving no graded, server-side notion of ``how risky is this particular call.''
Second, general vulnerability scores such as CVSS are demonstrably insensitive to
AI-specific threats. Third, even the systems closest to ours ---
AgentBound~\cite{buhler2026agentbound} and PAuth~\cite{sharma2026pauth} --- issue
an allow/deny decision rather than a graduated risk score. The proposed research
addresses all three by placing an interpretable, factor-based risk scorer at the
server boundary, computed statically from a tool's declared properties and
dynamically from the specific request.
